\chapter{Sprint 3 --- Feature B: Quản lý đội xe (\texttt{wujia\_fleet} + \texttt{wujia\_delivery})}

\section{Bối cảnh nghiệp vụ}

Sau khi xong mục A (Quản lý nhượng quyền), Wujia cần số hóa luồng \emph{điều phối giao hàng} từ kho công ty xuống các cửa hàng nhượng quyền. Hiện tại đang quản lý qua Excel + Zalo, không tracking được:

\begin{itemize}
    \item Mỗi tháng tốn bao nhiêu tiền cước cho nhà xe Nguyễn Dũng vs xe nội bộ?
    \item Cửa hàng HCM-01 tốn cước nhiều hơn HN-01 không? (để đối soát chia phí với chủ tiệm).
    \item Xe nào hay vượt tải, xe nào ít dùng?
    \item Bảng giá đi tỉnh năm nay tăng bao nhiêu so với năm ngoái?
\end{itemize}

Feature B trả lời các câu hỏi này bằng 2 module backend: \texttt{wujia\_fleet} (đăng ký xe + bảng giá) và \texttt{wujia\_delivery} (luồng điều phối + báo cáo cước). Portal cho cửa hàng theo dõi giao hàng được defer (per ADR-009).

\section{Quyết định kiến trúc (2026-05-02)}

Gộp 3 module BA đề xuất (\texttt{wujia\_fleet}, \texttt{wujia\_delivery}, \texttt{wujia\_delivery\_cost}) thành 2 module:

\begin{itemize}
    \item \texttt{wujia\_fleet} --- master data: nhà xe / loại xe / xe / bảng giá. Depends \texttt{wujia\_core}, \texttt{wujia\_franchise}.
    \item \texttt{wujia\_delivery} --- extension \texttt{stock.picking} + \texttt{stock.picking.batch} + cost compute + 2 SQL view report. Depends \texttt{wujia\_fleet}, \texttt{wujia\_sale}, \texttt{stock\_picking\_batch}.
\end{itemize}

Lý do: tránh module quá nhỏ (delivery\_cost), tránh circular dep cost vs batch. Field \texttt{current\_batch\_id} trên vehicle đặt trong \texttt{wujia\_delivery} (không phải \texttt{wujia\_fleet}) vì nó cần biết \texttt{stock.picking.batch.vehicle\_id} --- do \texttt{wujia\_delivery} thêm vào.

\section{Module \texttt{wujia\_fleet} --- Sổ tay đội xe}

\subsection{Nhà xe (Fleet Provider) --- "Ai cho mình thuê?"}

Wujia có 2 loại đội xe:

\begin{itemize}
    \item \textbf{Xe công ty}: của chính Wujia, tài xế nhân viên.
    \item \textbf{Xe thuê ngoài}: gọi từ công ty vận tải bên ngoài (vd CTY Nguyễn Dũng) --- mỗi tháng phải xuất Purchase Order trả tiền.
\end{itemize}

Mỗi nhà xe ghi: tên, mã, người liên hệ, SĐT, email. Nếu thuê ngoài --- link với 1 record Vendor (\texttt{res.partner}) trong Odoo. Có chatter tracking thay đổi.

\subsection{Loại xe (Fleet Type) --- "Xe to nhỏ ra sao?"}

Định nghĩa các "model" xe Wujia hay dùng: Bán tải 0.9T, Tải 1.9T, Tải 3.5T, Đông lạnh 17T... Field quan trọng: \texttt{payload\_capacity\_ton} (tải trọng tối đa, để hệ thống cảnh báo vượt tải sau này).

\subsection{Xe (Vehicle) --- "Cụ thể từng chiếc"}

Từng chiếc xe ngoài đời. Ghi biển số, mã nội bộ, tài xế chính, SĐT. Quan trọng nhất: \textbf{trạng thái xe} 6-state (Sẵn sàng / Vào bãi / Đã sắp chuyến / Đang giao / Bảo trì / Ngưng dùng). Có kanban view group theo trạng thái --- nhìn 1 phát biết "5 xe sẵn sàng, 2 xe đang giao".

Có thể gắn QR code (in dán lên xe) cho check-in tại bãi.

\subsection{Bảng giá vận chuyển (Pricelist) --- "Cước phí thế nào?"}

Cốt lõi tài chính. Cước KHÁC NHAU theo: loại xe + khu vực giao + nhà xe + thời điểm. Nên 1 bảng giá ghi:

\begin{quote}
\textit{"Bảng giá xe Tải 3.5T của Nguyễn Dũng năm 2026, gồm các dòng:}
\begin{itemize}
    \item \textit{HCM Q1 → 800k cước + 50k drop}
    \item \textit{Hà Nội → 2.5M cước + 100k drop}
    \item \textit{Đà Nẵng → 1.5M + 75k drop"}
\end{itemize}
\end{quote}

Mỗi dòng = mức cước cho 1 nhóm khu vực + phí drop (phí mỗi điểm dừng giao). \textbf{Vòng đời} 4 trạng thái: Nháp → Đang áp dụng → Hết hạn (auto qua cron daily khi quá ngày) → Lưu trữ.

\subsection{Phân quyền}

2 group trong privilege "Wujia Fleet":
\begin{itemize}
    \item \texttt{group\_fleet\_user} --- xem master data, không sửa.
    \item \texttt{group\_fleet\_manager} --- full CRUD.
\end{itemize}

\subsection{Menu Fleet Management}

\begin{lstlisting}[style=shell]
Fleet Management
+-- Operation
|   +-- Vehicles            (xe)
|   +-- Pricelists          (bang gia)
|   +-- Delivery Orders     (phieu xuat - tu wujia_delivery)
|   +-- Batch Management    (chuyen xe - tu wujia_delivery)
+-- Reports                 (tu wujia_delivery)
|   +-- Fleet Expense Report
|   +-- Franchise Shipping Cost Report
+-- Configuration           (group_fleet_manager)
    +-- Fleet Providers     (nha xe)
    +-- Fleet Types         (loai xe)
\end{lstlisting}

\section{Module \texttt{wujia\_delivery} --- Quy trình điều phối}

Luồng nghiệp vụ \textbf{end-to-end}:

\subsection{Bước 1: Cửa hàng đặt hàng (đã có từ wujia\_sale)}

Cô Anh chủ tiệm HN-01 vào portal đặt 100kg trà sữa nguyên liệu. Hệ thống tạo Sale Order.

\subsection{Bước 2: Confirm SO → tự sinh phiếu xuất}

Khi admin/auto confirm SO, Odoo tự tạo phiếu xuất kho \texttt{WH/OUT/00123}. \textbf{Module wujia\_delivery vào việc}: override \texttt{sale.order.\_action\_confirm} → copy \texttt{franchise\_id = HN-01} từ SO sang picking. Khu vực (\texttt{area\_id}) tự lấy theo cửa hàng (related store).

→ Mục đích: phòng điều xe cần biết picking này giao cho cửa hàng nào, ở khu vực nào, để chọn xe phù hợp.

\subsection{Bước 3: Phòng điều xe gom phiếu thành chuyến (Batch)}

Cuối ngày, phòng điều xe nhìn 50 phiếu, gom theo tuyến: "5 phiếu HN-01, HN-02, HN-03 chung 1 chuyến đi Hà Nội". → Tạo 1 Batch mới, gán 5 picking.

\subsection{Bước 4: Chọn xe cho chuyến}

Chọn xe \texttt{51C-34567 (Tải 3.5T)}. Hệ thống ngay lập tức:

\begin{itemize}
    \item Tự điền nhà xe + loại xe + tải trọng (related store, đỡ gõ tay).
    \item Tính tổng kg các phiếu trong chuyến (vd 600kg).
    \item Hiển thị \% sử dụng tải = 17.14\% (600/3500).
    \item Nếu vượt 3500kg → \textbf{badge đỏ "Vượt tải"} + ghi rõ thiếu bao nhiêu kg.
\end{itemize}

\subsection{Bước 5: Tự đề xuất bảng giá + tính cước}

Hệ thống search bảng giá phù hợp (Tải 3.5T + Nguyễn Dũng + ngày hôm nay) → \textbf{tự đề xuất}. Người dùng có thể đổi tay.

Tính cước:
\begin{enumerate}
    \item Nhìn các khu vực batch giao (HN-01, HN-02, HN-03 đều thuộc Hà Nội).
    \item Lấy giá CAO NHẤT trong các dòng giá match khu vực → cước cơ bản (vd 2.5M).
    \item Cộng phí drop cho từng điểm (vd 100k × 3 = 300k).
    \item \textbf{Tổng = 2.8M}.
\end{enumerate}

\subsection{Bước 6: Phân bổ ngược cước về từng phiếu}

Để báo cáo "cửa hàng nào tốn nhiều cước nhất", hệ thống chia cước theo \textbf{tỉ lệ kg}:

\begin{lstlisting}[style=shell]
HN-01 (200kg / 600kg = 33%) -> 833k cuoc + 100k drop = 933k
HN-02 (200kg / 600kg = 33%) -> 933k
HN-03 (200kg / 600kg = 33%) -> 933k
\end{lstlisting}

Mỗi phiếu xuất sẽ có 2 cột: \texttt{shipping\_cost}, \texttt{drop\_fee} (đã phân bổ).

\subsection{Bước 7: Theo dõi trạng thái chuyến}

Phòng điều xe bấm các nút:

\begin{enumerate}
    \item "Đã gán xe" (Draft → Assigned) sau khi xác nhận xe.
    \item "Bắt đầu chất hàng" (Assigned → Loading).
    \item "Xuất phát" (Loading → Delivering) → tự ghi \texttt{actual\_departure = now()}, đẩy tất cả phiếu trong chuyến → "Đang giao".
    \item "Hoàn tất giao" (Delivering → Done).
    \item "Hủy chuyến" --- nếu sự cố.
\end{enumerate}

Trạng thái này CHẠY SONG SONG với state chuẩn của Odoo (\texttt{draft}/\texttt{in\_progress}/\texttt{done}) --- không thay thế.

\subsection{Bước 8: Báo cáo}

\textbf{Fleet Expense Report} (chi phí đội xe): mỗi dòng = 1 chuyến. Pivot theo nhà xe + tháng → biết "T5 thuê Nguyễn Dũng tốn bao nhiêu, xe nội bộ tốn bao nhiêu". Filter "Vượt tải" → liệt kê các chuyến từng vượt tải.

\textbf{Franchise Shipping Cost Report} (chi phí theo cửa hàng): mỗi dòng = 1 phiếu xuất. Pivot theo cửa hàng + tháng → biết "HN-01 tháng này tốn 5M cước, HCM-01 tốn 3M". Dùng để \textbf{đối soát} với hợp đồng nhượng quyền.

\section{Performance optimizations (2026-05-02 patch)}

Sau khi triển khai cốt lõi, đã thực hiện các tối ưu sau để ready cho scale tương lai (5000+ batch/tháng, 50000+ picking):

\subsection{17 BTREE index thêm trên hot fields}

\begin{longtable}{p{6cm} p{8cm}}
\toprule
\textbf{Field} & \textbf{Lý do} \\
\midrule
\endhead
\texttt{wujia\_fleet\_provider.provider\_type} & Filter "Xe công ty/Thuê ngoài" \\
\texttt{wujia\_fleet\_management.provider\_id} & FK + filter theo nhà xe \\
\texttt{wujia\_fleet\_management.fleet\_type\_id} & FK + group by loại xe \\
\texttt{wujia\_fleet\_management.license\_plate} & Search nhanh theo biển số \\
\texttt{wujia\_fleet\_management.vehicle\_status} & Kanban group + filter "Sẵn sàng" \\
\texttt{wujia\_fleet\_pricelist.fleet\_type\_id} & Auto-suggest pricelist theo loại xe \\
\texttt{wujia\_fleet\_pricelist.provider\_id} & Auto-suggest pricelist theo nhà xe \\
\texttt{wujia\_fleet\_pricelist.state} & Cron expire + filter "Đang áp dụng" \\
\texttt{wujia\_fleet\_pricelist.date\_from} & Auto-suggest pricelist theo ngày \\
\texttt{wujia\_fleet\_pricelist.date\_to} & Cron expire WHERE date\_to < today \\
\texttt{stock\_picking.vehicle\_id} & Filter list picking theo xe \\
\texttt{stock\_picking.provider\_id} & Filter list picking theo nhà xe \\
\texttt{stock\_picking.delivery\_status} & Filter "Chờ điều phối / Đang giao" \\
\texttt{stock\_picking\_batch.vehicle\_id} & FK + lookup current\_batch \\
\texttt{stock\_picking\_batch.provider\_id} & Group report by nhà xe \\
\texttt{stock\_picking\_batch.fleet\_type\_id} & Group report by loại xe \\
\texttt{stock\_picking\_batch.pricelist\_id} & FK \\
\texttt{stock\_picking\_batch.delivery\_batch\_status} & Filter "Đang giao" \\
\bottomrule
\end{longtable}

\subsection{Store \texttt{vehicle\_count}}

Trước: compute mỗi lần read list provider → O(n) query (1 per provider). Sau: \texttt{store=True, compute\_sudo=True} → 1 query duy nhất, auto-recompute khi vehicle.active đổi.

\subsection{Tối ưu \texttt{\_compute\_pricelist\_id}}

Trước: mỗi batch trong list view trigger 1 SQL search pricelist → 100 batch = 100 query. Sau: prefetch tất cả pricelist active 1 lần, filter in-memory → 1 query duy nhất cho 100 batch.

\section{Cách verify}

\subsection{Drop + reinstall fresh}

\begin{lstlisting}[style=shell]
PGPASSWORD=1 dropdb -h 127.0.0.1 -U odoo19 wujia_tea_19
PGPASSWORD=1 createdb -h 127.0.0.1 -U odoo19 wujia_tea_19
cd /home/huyban/odoo-dev/WujiaTea/odoo19
/home/huyban/miniconda3/envs/odoo/bin/python odoo-bin \\
  -c /home/huyban/odoo-dev/WujiaTea/config/odoo.conf -d wujia_tea_19 \\
  -i wujia_core,wujia_franchise,wujia_sale,wujia_portal_base, \\
     wujia_portal_layout,wujia_fleet,wujia_delivery \\
  --stop-after-init --without-demo=False
\end{lstlisting}

\subsection{Schema check}

\begin{lstlisting}[style=shell]
psql -d wujia_tea_19 -c "\\d wujia_fleet_management"
psql -d wujia_tea_19 -c "\\d stock_picking" | grep -E "(franchise|vehicle|delivery)"
psql -d wujia_tea_19 -c "SELECT viewname FROM pg_views WHERE viewname LIKE 'wujia_%'"
\end{lstlisting}

Kết quả:
\begin{itemize}
    \item 5 table mới (\texttt{wujia\_fleet\_provider/type/management/pricelist/pricelist\_line}).
    \item \texttt{stock\_picking}: thêm 9 cột; \texttt{stock\_picking\_batch}: thêm 14 cột.
    \item 2 SQL view: \texttt{wujia\_fleet\_shipping\_report}, \texttt{wujia\_franchise\_shipping\_report}.
    \item 17 BTREE index mới.
\end{itemize}

\subsection{End-to-end test cost compute}

Script test ở \texttt{/tmp/test\_cost\_e2e.py}: tạo 2 picking gắn franchise + area, gom thành batch + chọn vehicle, verify auto-suggest pricelist + cost compute + allocation.

Kết quả thực tế (đã pass):

\begin{lstlisting}[style=shell]
Vehicle: 51C-34567 -- Tai 3.5T (Nguyen Dung) cap=3500kg
Pricelist: PL-T35-2026 (HCM 800k + 50k drop, HN 2.5M + 100k drop)

Picking p1: HCM-01, planned_weight=200kg
Picking p2: HN-01,  planned_weight=400kg

Batch BATCH/OUT/00003:
  planned_weight = 600 kg
  capacity_usage = 17.14% (600/3500)
  is_over_capacity = False
  pricelist (auto-suggest) = PL-T35-2026
  shipping_cost   = 2,500,000 (max gia trong cac line match area)
  drop_fee_total  =   150,000 (50k HCM + 100k HN)
  total_cost      = 2,650,000

Phan bo nguoc:
  p1 (HCM, 200kg/600kg = 33%): shipping=833,333  drop=50,000
  p2 (HN,  400kg/600kg = 67%): shipping=1,666,667 drop=100,000

Action_delivery_start():
  batch.delivery_batch_status = delivering
  batch.actual_departure = 2026-05-02 13:05:07
  p1.delivery_status = delivering, p2.delivery_status = delivering

Reports:
  Fleet report: planned=600 total=2,650,000
  Franchise report:
    HCM-01: weight=200 total_alloc=883,333
    HN-01:  weight=400 total_alloc=1,766,667
\end{lstlisting}

\section{Demo data --- không load qua manifest}

BA quyết định: KHÔNG đặt demo data trong file XML để tránh đẩy lên server. Script local:

\begin{lstlisting}[style=shell]
cd /home/huyban/odoo-dev/WujiaTea/odoo19
/home/huyban/miniconda3/envs/odoo/bin/python odoo-bin shell \\
  -c /home/huyban/odoo-dev/WujiaTea/config/odoo.conf \\
  -d wujia_tea_19 --no-http \\
  < /home/huyban/odoo-dev/WujiaTea/scripts/seed_fleet_demo.py
\end{lstlisting}

Idempotent (lookup theo \texttt{code}). Tạo: 2 provider + 4 type + 5 vehicle + 1 pricelist 2 line.

% ===========================================================================
